% !TEX program = xelatex
\documentclass[10pt, a4paper]{article}
% ===== Essential Packages =====
\usepackage[utf8]{inputenc}       % UTF-8 encoding
\usepackage[T1]{fontenc}          % Better font encoding
\usepackage{lmodern}              % Modern Latin Modern font
\usepackage{ctex}                 % Chinese support (must be loaded after fontenc)
\usepackage{subcaption}
% ===== Math & Symbols =====
\usepackage{amsmath}              % Advanced math
\usepackage{amssymb}              % Additional math symbols
\usepackage{mathrsfs}             % Ralph Smith's Formal Script
\usepackage{braket}               % Dirac bra-ket notation

% ===== Graphics & Tables =====
\usepackage{graphicx}             % Image support
\usepackage{float}                % Better float control
\usepackage{wrapfig}              % Wrapped figures
\usepackage{tikz}                 % Vector graphics
\usepackage[table]{xcolor}        % Colored tables
\usepackage{array}                % Enhanced tables
\usepackage{multirow}             % Multirow tables
\usepackage{tabularx}             % Adjustable-width tables
\usepackage{booktabs}             % Professional table lines
\usepackage{siunitx}              % SI units in tables
\usepackage{caption}              % Better captions
\usepackage{adjustbox}            % Box resizing

% ===== Code Listings =====
\usepackage{listings}             % Code blocks
\usepackage{tcolorbox}            % Colored boxes
\tcbuselibrary{listings, breakable} % Listings in tcolorbox

% ===== Misc =====
\usepackage{lipsum}               % Dummy text
\usepackage{footmisc}             % Better footnotes
\usepackage[margin=1in]{geometry} % Page margins
\usepackage{diagbox}
\usepackage{ctex}
\author{张世航}
\title{Zeeman效应}
\begin{document}
	\maketitle
	\section{荷质比的计算}

	\begin{tabular}{|c|c|c|c|}
		\hline
		B & 0 & 0 & 0\\
		\hline
		\(r^1\)&1.16mm&1.17mm&1.16mm\\
		\hline
		\(r^1\)&1.45mm&1.44mm&1.44mm\\
		\hline
		
	
		
		\hline
		B & 1052mT & 1054mT & 1059mT\\
		\hline
		\(r^1_{\text{内}}\)&1.12mm&1.12mm&1.12mm\\
		\hline
		\(r^1_{\text{中}}\)&1.18mm&1.18mm&1.18mm\\
		\hline
		\(r^1_{\text{外}}\)&1.23mm&1.23mm&1.23mm\\
		\hline
		\(r^2_{\text{内}}\)&1.40mm&1.40mm&1.40mm\\
		\hline
		\(r^2_{\text{中}}\)&1.45mm&1.45mm&1.46mm\\
		\hline
		\(r^2_{\text{外}}\)&1.51mm&1.51mm&1.51mm\\
		\hline
	\end{tabular}\\
	我们考虑两束具有微小波长差的单程光\(\lambda_1,\lambda_2\),其中\(\lambda_1>\lambda,\lambda_1\approx\lambda_2\approx\lambda\),
	如果\(\lambda_1\)和\(\lambda_2\)的波长差逐渐增大,使得\(\lambda_2\)的k序与\(\lambda_1\)的\(k-1\)序花纹重合,这时以下条件得以满足
	\begin{equation}
		k\lambda_2 = (k-1)\lambda_1
	\end{equation}
	\begin{equation}
		\Delta \lambda = \lambda_1-\lambda_2 = \lambda_1-\frac{k-1}{k}\lambda_1=\frac{\lambda_1}{k}\approx\frac{\lambda^2}{2d}
	\end{equation}
	用波数表示为
	\begin{equation}
		\Delta \tilde{\nu} = \frac{1}{2d}
	\end{equation}
	对于F-P干涉仪
	\begin{equation}
		\cos\theta = \frac{f}{\sqrt{f^2+(\frac{D}{2})^2}}\approx1-\frac{1}{8}\frac{D^2}{f^2}
	\end{equation}
	对于衍射条件
	\begin{equation}
		2d\cos\theta = k\lambda
	\end{equation}
	则有
	\begin{equation}
		2d(1-\frac{1}{8}\frac{D^2_k}{f^2}) = k\lambda
	\end{equation}
	可以得到
	\begin{equation}
		D^2_k =8f^2(1 - \frac{k\lambda}{2d} )
	\end{equation}
	从而条纹间距为
	\begin{equation}
		\Delta D^2 = D_{k-1}^2 - D_k^2 = \frac{4f^2\lambda}{d}
	\end{equation}
	我们也可以测出同一序中不同波长\(\lambda_a,\lambda_b\)之差
	\begin{equation}
	   \lambda_a-\lambda_b = \frac{d}{4f^2k}(D_b^2-D_a^2) = \frac{\lambda}{k}\frac{D_b^2-D_a^2}{D_{k-1}^2-D_k^2}
	\end{equation}
	用波数表示为
	\begin{equation}
		\Delta\tilde{\nu} = \tilde{\nu}_b - \tilde{\nu}_a = \frac{1}{2d}\frac{\Delta D_{ab}^2}{\Delta D^2}
	\end{equation}
	
	我们可以计算得到\\
	\begin{tabular}{|c|c|c|c|c|c|c|c|c|c|c|c|c|}
		\hline
		n & 1 & 2 & 3&4&5&6&7&8&9&10&11&12\\
		\hline
		\(\frac{\Delta D_{ab}^2}{mm^2}\)&0.138&  0.138&  0.138&  0.1205 &0.1205& 0.1205 &0.1425& 0.1425 &0.1716& 0.1475&0.1776& 0.1485\\
		\hline
        \(\frac{\Delta D^2}{mm^2}\)&2.8325& 2.8325 &2.8896& 2.8944 &2.9535 &3.0128& 2.8325 &2.8325 &2.8896& 2.8944&2.9535 &3.0128\\
		\hline
	\end{tabular}\\
	根据我们的数据
	\begin{equation}
		\Delta D^2 = D_{k-1}^2 -D_k^2= 2.8791111111111096mm^2
	\end{equation}
	可以得到
	\begin{equation}
		\Delta\bar{ D^2}_{ab} = 0.14214166666666667 mm
	\end{equation}
	则
	\begin{equation}
    \Delta \bar{\tilde{\nu}} = \frac{1}{2d} \frac{\Delta\bar{D^2}_{ab}}{\Delta D^2} = 24.684991509725236\,\text{m}^{-1}
    \end{equation}
	\begin{equation}
		\frac{e}{m} = \frac{4\pi c}{\frac{B}{2}}\Delta \bar{\tilde{\nu}}=176417489245.60202C/kg
	\end{equation}
	\section{左右旋光判断}
	在\(\theta = \frac{\pi}{4},\frac{5\pi}{4}\)的时候出现外圈,在\(\theta = \frac{3\pi}{4},\frac{7\pi}{4}\)的时候出现内圈.\\
	\begin{tabular}{|c|c|c|c|}
		\hline
		内圈&右旋光&角动量为-h&\(\Delta E\)下降\\
		\hline
        外圈&左旋光&角动量为h&\(\Delta E\)升高\\
		\hline
	\end{tabular}\\
\end{document} 